\lecture{8.9}{Longitudinal TMLE}
\label{lec:8-9}

\subsection{Sequential TMLE Algorithm}

Targeted minimum loss-based estimation (TMLE) modifies sequential
g-computation so that the final estimator remains a plug-in estimator while also
solving bias-correction equations related to the efficient influence function.
In longitudinal settings, TMLE proceeds backward through time, just like
sequential regression, but inserts a targeting step after each outcome
regression.

For the always-treated regime, the algorithm has the following structure:
\begin{enumerate}
  \item Estimate the final outcome regression
  $\Qbar_{K+1,n}(\Lbar_K)=\E_n(Y\given \Abar_K=\mathbf{1},\Lbar_K)$.
  \item Target this regression using a clever covariate based on
  $1/\gbar_{K,n}$, producing $\Qbar_{K+1,n}^*$.
  \item Regress $\Qbar_{K+1,n}^*$ on earlier history among individuals who
  followed the regime through time $K-1$, producing $\Qbar_{K,n}$.
  \item Target $\Qbar_{K,n}$ using the clever covariate based on
  $1/\gbar_{K-1,n}$, producing $\Qbar_{K,n}^*$.
  \item Continue backward until $\Qbar_{1,n}^*$ is obtained.
  \item Estimate the target mean by
  \[
    \psi_{n,\TMLE}=\frac{1}{n}\sum_{i=1}^n \Qbar_{1,n}^*(L_{0i}).
  \]
\end{enumerate}

\subsection{Targeted Refinement at Each Time}

Assume the outcome or pseudo-outcome is bounded in $[0,1]$, so logistic
fluctuation is natural. At a generic step $j$, suppose we have an initial
regression $\Qbar_{j,n}$ and a pseudo-outcome $Z_j$. The targeting step fits a
low-dimensional regression with:
\begin{itemize}
  \item outcome $Z_j$,
  \item offset $\logit\{\Qbar_{j,n}\}$,
  \item clever covariate based on the inverse probability of following the
  regime through the relevant time.
\end{itemize}
The fitted fluctuation coefficient updates $\Qbar_{j,n}$ to
$\Qbar_{j,n}^*$.

\subsection{Clever Covariates}

For two treatment times and regime $(1,1)$, the final outcome regression is
\[
  \Qbar_{2,n}(L_0,L_1)\approx \E(Y\given A_0=1,A_1=1,L_0,L_1).
\]
The clever covariate for targeting this regression is
\[
  H_{2,n}=\frac{\ind(A_0=1,A_1=1)}
  {g_{0,n}(L_0)g_{1,n}(\Lbar_1)}.
\]
A logistic fluctuation fits
\[
  \logit\{\Qbar_{2,n}^*(L_0,L_1)\}
  =
  \logit\{\Qbar_{2,n}(L_0,L_1)\}
  +\hat\epsilon_2 H_{2,n}.
\]
The next regression estimates
\[
  \Qbar_{1,n}(L_0)\approx
  \E\{\Qbar_{2,n}^*(L_0,L_1)\given A_0=1,L_0\},
\]
and its clever covariate is
\[
  H_{1,n}=\frac{\ind(A_0=1)}{g_{0,n}(L_0)}.
\]
The update is
\[
  \logit\{\Qbar_{1,n}^*(L_0)\}
  =
  \logit\{\Qbar_{1,n}(L_0)\}+\hat\epsilon_1H_{1,n}.
\]

\subsection{Final Plug-In Estimate}

The final estimate is
\[
  \psi_{n,\TMLE}(1,1)=\frac{1}{n}\sum_{i=1}^n\Qbar_{1,n}^*(L_{0i}).
\]
It is a plug-in estimator because it is the empirical mean of an updated
counterfactual outcome regression. For binary outcomes, this guarantees that the
estimate lies in $[0,1]$, unlike some weighted estimators that can behave poorly
in small samples.

\subsection{TMLE Implementation in R}

The following sketch shows the two-timepoint binary-outcome implementation.

\begin{lstlisting}[language=R]
# propensity scores
ps0_fit <- glm(A0 ~ L0, family = binomial())
ps1_fit <- glm(A1 ~ L0 + A0 + L1, family = binomial())

g_0 <- predict(ps0_fit, type = "response")
g_1 <- predict(
  ps1_fit,
  newdata = data.frame(L0 = L0, A0 = 1, L1 = L1),
  type = "response"
)

# initial Q2
or2_fit <- glm(Y ~ A0 + A1 + L0 + L1, family = binomial())
Q2_obs <- predict(or2_fit, type = "response")
Q2_11 <- predict(
  or2_fit,
  newdata = data.frame(A0 = 1, A1 = 1, L0 = L0, L1 = L1),
  type = "response"
)

# target Q2
H2 <- (A0 == 1 & A1 == 1) / (g_0 * g_1)
eps2_fit <- glm(
  Y ~ -1 + offset(qlogis(Q2_obs)) + H2,
  family = binomial()
)
Q2_11_star <- plogis(qlogis(Q2_11) + coef(eps2_fit)["H2"] / (g_0 * g_1))

# initial Q1
or1_fit <- glm(Q2_11_star ~ A0 * L0, family = binomial())
Q1_obs <- predict(or1_fit, type = "response")
Q1_1 <- predict(
  or1_fit,
  newdata = data.frame(A0 = 1, L0 = L0),
  type = "response"
)

# target Q1
H1 <- (A0 == 1) / g_0
eps1_fit <- glm(
  Q2_11_star ~ -1 + offset(qlogis(Q1_obs)) + H1,
  family = binomial()
)
Q1_1_star <- plogis(qlogis(Q1_1) + coef(eps1_fit)["H1"] / g_0)

psi_tmle <- mean(Q1_1_star)
\end{lstlisting}

\subsection{Relationship to Sequential Regression}

Sequential regression estimates the nested g-computation formula and stops.
Longitudinal TMLE follows the same backward structure, but after each regression
it performs a targeted update using the propensity score sequence. This update
is designed so that the empirical mean of the efficient influence-function
component is close to zero.

\begin{keyidea}
Longitudinal TMLE is targeted sequential regression. It keeps the plug-in
interpretation of g-computation while incorporating inverse-probability
information through clever covariates.
\end{keyidea}
